\subsection{Reading a polar equation as a radial instruction}

A polar equation
\[
r=f(\theta)
\]
should be read operationally: choose a direction $\theta$, then travel the
prescribed distance $r$ from the pole.

The simplest example is
\[
\boxed{r=R},
\]
which gives the circle of radius $R$ centred at the pole.

A more revealing example is
\[
\boxed{r=2R\cos\theta}.
\]
Using $x=r\cos\theta$ and $r^2=x^2+y^2$,
\[
r^2=2Rr\cos\theta
\quad\Longrightarrow\quad
x^2+y^2=2Rx,
\]
so
\[
\boxed{(x-R)^2+y^2=R^2}.
\]
The same curve is therefore a circle centred at $(R,0)$ and passing through the
pole.

\begin{center}
\begin{tikzpicture}[scale=0.9,>=stealth]
    \draw[axis,->] (-1.8,0)--(6.3,0) node[right] {$x$};
    \draw[axis,->] (0,-3.3)--(0,3.5) node[above] {$y$};
    \draw[subset] (2.0,0) circle (2.0);
    \coordinate (P) at ({4*cos(45)*cos(45)},{4*cos(45)*sin(45)});
    \draw[blue!70!black,line width=1.1pt,->] (0,0)--(P)
        node[midway,above left] {$r=2R\cos\theta$};
    \draw (0.9,0) arc[start angle=0,end angle=45,radius=0.9];
    \node at (1.15,0.45) {$\theta$};
    \fill[geometricSubsetColor] (0,0) circle (2.2pt) node[below left] {pole};
    \fill[geometricSubsetColor] (2.0,0) circle (2.2pt) node[below] {centre};
    \fill[geometricSubsetColor] (P) circle (2.2pt) node[above right] {$P$};
\end{tikzpicture}
\end{center}

\begin{interpretationbox}[What the equation is doing]
The factor $\cos\theta$ shortens the permitted radial distance as the ray turns
away from the positive axis. At $\theta=\pi/2$ the distance becomes zero, so the
curve returns to the pole. A Cartesian equation hides this directional story;
the polar equation displays it directly.
\end{interpretationbox}

This way of reading a formula prepares the conic equation. The next question is
whether one radial law can generate ellipse, parabola, and hyperbola by changing
only one parameter.
